\section{Data Release Processing} \label{sec:drp}

Data Release Processing (DRP) is the data pipeline that produces the calibrated images, detection catalogs, and derived data 
products defined in section \S TBD, using the LSST Science Pipelines (citation). DP2 was processed at the United States Data Facility (USDF) at SLAC, with pilot runs at the National Institute of Nuclear and Particle Physics (IN2P3) computing center. 

The processing campaign consists of four major stages: (1) single-frame processing, (2) calibration, (3) multi-visit processing, (4) variable and transient source processing.
Data quality and key metrics for verification and validation are actively monitored during each stage, and comprehensive verification and validation assessments are completed before processing the next stage.

This section describes the high-level algorithms, configurations, and verification metrics for each stage, including changes and improvements from DP1 (citation).

\subsection{Single-Visit Processing} \label{sec:drp-singlevisit}

\subsubsection{Calibration \& Instrument Signature Removal (ISR)} \label{sec:drp-singlevisit-calibrationisr}
\subsubsection{Background Subtraction} \label{sec:drp-singlevisit-background}

\subsection{Calibration} \label{sec:drp-calibration}
\subsubsection{PSF Modeling} \label{sec:drp-calibration-psfmodeling}

PSF modeling in DP2 uses two distinct approaches, as described in the 
\citep{PSTN-019} and it is similar to what was originally 
described in \citealt{DP1}. The first model—faster but less precise—is fit 
with PSFEx \citep{2011ASPC..442..435B}, while the final model—slower yet more 
accurate—is fit with Piff \citep{DES:2020vau}. Broadly speaking, PSFEx 
and Piff both represent the PSF on a pixel basis and interpolate its 
parameters across a single CCD using two-dimensional polynomials.
However, their technical implementations differ (see Bertin et al. 
and \citealt{DES:2020vau} for more details). Previous comparisons have shown 
that Piff outperforms PSFEx by delivering an unbiased reconstruction of 
the PSF size (see \citealt{2018MNRAS.481.1149Z} and \citealt{DES:2020vau}), can work in 
sky coordinates like in \citealt{2025OJAp....8E..26S}, can include a color dependency like 
in \citealt{2025OJAp....8E..26S}, and is supporting alternative models and interpolation 
schemes (see \citealt{2025OJAp....8E..26S} for more details).

In the initial PSF modeling pass, PSFEx is run with a pixel basis and 
a second-order polynomial interpolation on each CCD. For the final modeling, 
Piff uses its PixelGrid with a fourth-order polynomial interpolation per 
CCD—except in the u-band, where star counts are too low to constrain a 
fourth-order fit, so a second-order polynomial is used instead. Both are 
working in pixel coordinates. Choices of polynomial order, overall PSF-modeling 
performance, and known issues are discussed in Section 
“Performance and Known Issues for PSF.”

\subsubsection{Astrometric Calibration} \label{sec:drp-calibration-astrometry}
\subsubsection{Photometric Calibration} \label{sec:drp-calibration-photometry}

\subsection{Multi-Visit Processing} \label{sec:drp-multivisit}
\subsubsection{Input Data Selection} \label{sec:drp-multivisit-inputselection}
\subsubsection{Coaddition} \label{sec:drp-multivisit-coaddition}
\subsubsection{Detection, Deblending, \& Measurement} \label{sec:drp-multivisit-detection}
\subsubsection{First Look} \label{sec:drp-multivisit-firstlook}

\subsection{Variable \& Transient Source Processing} \label{sec:drp-variabletransient}
\subsubsection{Light Curves} \label{sec:drp-variabletransient-lightcurves}
\subsubsection{Solar System Processing} \label{sec:drp-variabletransient-solar}
\subsubsection{Difference Image Analysis (DIA)} \label{sec:drp-variabletransient-dia}
\subsubsection{Artifact Rejection} \label{sec:drp-variabletransient-artifact}
